%% Barrier Layer Depth -- Beamer presentation
%% Compile: pdflatex barrier_layers_beamer.tex
\documentclass[aspectratio=169, 12pt]{beamer}

%% ── Encoding & language ───────────────────────────────────────────────────────
\usepackage[utf8]{inputenc}
\usepackage[T1]{fontenc}
\usepackage[spanish]{babel}

%% ── Packages ──────────────────────────────────────────────────────────────────
\usepackage{graphicx}
\usepackage{booktabs}
\usepackage{array}
\usepackage{xcolor}
\usepackage{amsmath}
\usepackage{tikz}
\usepackage{colortbl}

%% ── Colors ────────────────────────────────────────────────────────────────────
\definecolor{oceanblue}{HTML}{065A82}
\definecolor{deepnavy}{HTML}{1C2951}
\definecolor{iceblue}{HTML}{CADCFC}
\definecolor{teal}{HTML}{1C7293}
\definecolor{lightgray}{HTML}{F4F4F4}

%% ── Theme ─────────────────────────────────────────────────────────────────────
\usetheme{default}
\setbeamertemplate{navigation symbols}{}
\setbeamertemplate{footline}{%
  \hfill\usebeamerfont{footline}\insertframenumber\,/\,\inserttotalframenumber\hspace{4pt}\vspace{4pt}%
}

\setbeamercolor{frametitle}{fg=white, bg=oceanblue}
\setbeamerfont{frametitle}{size=\large, series=\bfseries}
\setbeamercolor{title}{fg=white}
\setbeamercolor{subtitle}{fg=iceblue}
\setbeamercolor{author}{fg=iceblue}
\setbeamercolor{date}{fg=iceblue}
\setbeamercolor{normal text}{fg=black}
\setbeamercolor{structure}{fg=oceanblue}
\setbeamercolor{block title}{fg=white, bg=oceanblue}
\setbeamercolor{block body}{bg=lightgray}
\setbeamertemplate{itemize item}{\color{oceanblue}\textbullet}
\setbeamertemplate{itemize subitem}{\color{teal}--}

%% ── Helper: section divider frame ─────────────────────────────────────────────
\newcommand{\sectionframe}[1]{%
  \begin{frame}[plain, noframenumbering]
    \begin{tikzpicture}[remember picture, overlay]
      \fill[deepnavy] (current page.south west) rectangle (current page.north east);
      \node[text=white, font=\LARGE\bfseries, align=center]
        at (current page.center) {#1};
    \end{tikzpicture}
  \end{frame}
}

%% ── Helper: full-slide image ──────────────────────────────────────────────────
\newcommand{\fullimg}[1]{%
  \includegraphics[width=\textwidth, height=0.73\textheight, keepaspectratio]{#1}%
}

%% ── Caption style ─────────────────────────────────────────────────────────────
\newcommand{\captext}[1]{%
  \vskip2pt
  {\footnotesize\color{gray} #1}%
}

%% ── Title info ────────────────────────────────────────────────────────────────
\title{Barrier Layer Depth: Análisis Global 1993--2024}
\subtitle{Sensibilidad al método, comparación LR vs HR,\\tendencias y variabilidad interanual}
\date{2026}

%% ═══════════════════════════════════════════════════════════════════════════════
\begin{document}

%% ── TITLE ─────────────────────────────────────────────────────────────────────
{
\setbeamercolor{background canvas}{bg=deepnavy}
\begin{frame}[plain]
  \vfill
  \centering
  {\color{white}\Large\textbf{Barrier Layer Depth: Análisis Global 1993--2024}}\\[10pt]
  {\color{iceblue}\normalsize Sensibilidad al método, comparación LR vs HR,\\
  tendencias y variabilidad interanual}\\[30pt]
  {\color{iceblue}\small 2026}
  \vfill
\end{frame}
}

%% ═══════════════════════════════════════════════════════════════════════════════
\sectionframe{1. Física de la Barrier Layer}

%% ── Qué es la Barrier Layer ──────────────────────────────────────────────────
\begin{frame}{¿Qué es la Barrier Layer?}
\begin{columns}[c]

\column{0.52\textwidth}
\centering
\begin{tikzpicture}[scale=1.05]
  %% ── ocean column (x: 0..2, y: 0..6) ─────────────────────
  \fill[cyan!10]      (0,5.4) rectangle (2.0,6.0); % surface
  \fill[cyan!25]      (0,3.6) rectangle (2.0,5.4); % MLD
  \fill[teal!45]      (0,2.0) rectangle (2.0,3.6); % BLD  ← resaltado
  \fill[blue!60]      (0,0.0) rectangle (2.0,2.0); % thermocline
  %% ── boundaries ───────────────────────────────────────────
  \draw[white, line width=1.2pt] (0,5.4) -- (2.0,5.4);
  \draw[white, dashed, line width=1pt] (0,3.6) -- (2.0,3.6);
  \draw[white, line width=1.2pt] (0,2.0) -- (2.0,2.0);
  \draw[gray!40, thin] (0,0) rectangle (2.0,6.0);
  %% ── depth brackets (left) ────────────────────────────────
  \draw[<->, thick, black!50]
    (-0.25,3.6) -- (-0.25,5.4) node[midway,left,font=\scriptsize]{MLD};
  \draw[<->, thick, black!40]
    (-0.25,2.0) -- (-0.25,5.4) node[midway,left=3pt,font=\scriptsize,xshift=-4pt]{ILD};
  \draw[<->, very thick, teal]
    (-0.75,2.0) -- (-0.75,3.6)
    node[midway,left,font=\small\bfseries,color=teal]{BLD};
  %% ── right labels ─────────────────────────────────────────
  \node[right,font=\small\itshape]             at (2.2,5.70){Superficie};
  \node[right,font=\small\bfseries]            at (2.2,4.60){Capa de Mezcla};
  \node[right,font=\scriptsize,color=gray]     at (2.2,4.30){isohalina e isotérmica};
  \node[right,font=\small\bfseries,color=teal] at (2.2,2.90){Barrier Layer};
  \node[right,font=\scriptsize,color=gray]     at (2.2,2.60){fresca y cálida};
  \node[right,font=\scriptsize,color=gray]     at (2.2,2.30){estratif.\ salina};
  \node[right,font=\small\bfseries]            at (2.2,1.00){Termoclina};
\end{tikzpicture}

\column{0.45\textwidth}
\begin{block}{Definición}
\[ \text{BLD} = \text{ILD} - \text{MLD} \]
\end{block}
\vspace{6pt}
{\small
\begin{itemize}
  \item \textbf{BLD $>$ 0}: Barrier Layer presente
  \item \textbf{BLD $=$ 0}: ILD coincide con MLD
\end{itemize}
\vspace{8pt}
La BLD \textbf{inhibe el intercambio de calor} entre la superficie y la termoclina, modulando la SST y el acoplamiento océano-atmósfera.
}
\end{columns}
\end{frame}

%% ── Regiones ─────────────────────────────────────────────────────────────────
\begin{frame}{Regiones de análisis}
\begin{columns}[T]
\column{0.60\textwidth}
\fullimg{../displays/regions_map.png}

\column{0.38\textwidth}
\vspace{8pt}
\textbf{Datos:} CMEMS 0.25° mensual\\1993--2024 (384 meses)\\[10pt]
\textbf{4 regiones clave:}
\begin{itemize}
  \item {\color{red}\textbf{Índico Tropical}}\\20°S--20°N, 40°E--100°E
  \item {\color{violet}\textbf{Bahía de Bengala}}\\10°N--25°N, 80°E--100°E
  \item {\color{orange}\textbf{W. Pacific Warm Pool}}\\10°S--10°N, 130°E--180°E
  \item {\color{teal}\textbf{Amazon Plume}}\\5°S--10°N, 55°W--30°W
\end{itemize}
\end{columns}
\end{frame}

%% ═══════════════════════════════════════════════════════════════════════════════
\sectionframe{2. Datos y Métodos}

%% ── Productos CMEMS ──────────────────────────────────────────────────────────
\begin{frame}{Productos CMEMS utilizados}
\begin{center}
\begin{tabular}{lllp{4.2cm}}
\toprule
\textbf{Dataset ID} & \textbf{Resolución} & \textbf{Período} & \textbf{Variables} \\
\midrule
\texttt{...0.25deg\_P1M-m} & 0.25° mensual (LR) & 1993--2025 & thetao\_glor, so\_glor, mlotst\_glor \\[4pt]
\texttt{...0.083deg\_P1M-m} & 0.083° mensual (HR) & 1993--2025 & thetao, so, mlotst \\
\bottomrule
\end{tabular}
\end{center}
\vspace{10pt}
\begin{itemize}
  \item \textbf{MLD}: usamos el MLD pre-computado de CMEMS (\texttt{mlotst\_glor}) -- no lo recalculamos
  \item \textbf{ILD}: calculado con 6 configuraciones propias sobre los perfiles de temperatura
  \item \textbf{Grillas verticales:} LR tiene 74 niveles (0.5--5698 m); HR tiene 32 niveles (0.5--541 m)\\
        En los primeros 300 m relevantes para la BLD: LR 34 niveles, HR 28 niveles
\end{itemize}
\end{frame}

%% ── Grillas verticales ───────────────────────────────────────────────────────
\begin{frame}{Grillas verticales: LR vs HR (0--300 m)}
\begin{center}
\fullimg{../displays/vertical_grid_comparison.png}
\end{center}
\captext{Espaciado entre niveles en función de la profundidad. LR tiene más niveles que HR en los primeros 300 m.}
\end{frame}

%% ── 6 métodos ────────────────────────────────────────────────────────────────
\begin{frame}{6 métodos de detección de ILD}
\begin{center}
{\small
\begin{tabular}{@{} l l p{5.8cm} p{2.8cm} @{}}
\toprule
\textbf{N} & \textbf{Familia} & \textbf{Definición} & \textbf{Umbral} \\
\midrule
1 & Gradiente & Primera $z$ donde $\partial T/\partial z \leq$ umbral (desde 5\,m) & $-0.015$\,°C/m \\[2pt]
2 & Gradiente & ídem & $-0.025$\,°C/m \\[2pt]
3 & Gradiente & ídem & $-0.1$\,°C/m \\[2pt]
4 & Diferencia & Primera $z$ donde $T(z) < T(10\,\text{m}) - \Delta T$ & $\Delta T = 0.2$\,°C \\[2pt]
5 & Diferencia & ídem & $\Delta T = 0.5$\,°C \\[2pt]
6 & Diferencia & ídem & $\Delta T = 0.8$\,°C \\
\bottomrule
\end{tabular}
}
\end{center}
\vspace{6pt}
\[ \text{BLD} = \text{ILD} - \text{MLD} \quad \text{para los 6 métodos} \]
\captext{Objetivo: evaluar cuánto cambian los resultados según el método elegido.}
\end{frame}

%% ═══════════════════════════════════════════════════════════════════════════════
\sectionframe{3. Estructura Vertical -- Perfiles Individuales}

%% ── 6 métodos en un perfil ───────────────────────────────────────────────────
\begin{frame}{Los 6 métodos sobre un mismo perfil}
\begin{center}
\fullimg{../displays/sensitivity_one_point_2023-12-30_lat-10.0_lon55.0.png}
\end{center}
\captext{Lat $-10^{\circ}$, lon $55^{\circ}$, Dic 2023. Cada método detecta una ILD diferente, ilustrando la sensibilidad al umbral.}
\end{frame}

%% ═══════════════════════════════════════════════════════════════════════════════
\sectionframe{4. Sensibilidad al Método de Detección}

%% ── Ensemble maps ────────────────────────────────────────────────────────────
\begin{frame}{Mapa ensemble: media, spread y acuerdo entre métodos}
\begin{center}
\fullimg{../displays/sensitivity_ensemble_maps.png}
\end{center}
\captext{Mayor spread en el W. Pacific Warm Pool y Bahía de Bengala. Mayor acuerdo en el Atlántico tropical.}
\end{frame}

%% ── 6 configs side by side ───────────────────────────────────────────────────
\begin{frame}{BLD media anual: los 6 métodos comparados}
\begin{center}
\fullimg{../displays/summary_sensitivity_comparison.png}
\end{center}
\captext{Métodos de diferencia detectan Barrier Layer en $\sim$50\% del océano; los de gradiente en 24--35\%.}
\end{frame}

%% ── Stats table ──────────────────────────────────────────────────────────────
\begin{frame}{Estadísticas globales (1993--2024)}
\begin{columns}[T]
\column{0.50\textwidth}
\fullimg{../displays/sensitivity_stats_table.png}

\column{0.48\textwidth}
\vspace{8pt}
\begin{itemize}
  \item BLD media global: 12.7 m (grad $-0.015$) a 40.1 m (diff 0.8°C) -- factor $\sim3\times$
  \item Cobertura: 24--35\% (gradiente) vs 48--50\% (diferencia)
  \item Amplitud estacional: 10--21 m según método
\end{itemize}
\end{columns}
\end{frame}

%% ── Global time series ───────────────────────────────────────────────────────
\begin{frame}{Serie temporal global de BLD -- los 6 métodos (1993--2024)}
\begin{center}
\fullimg{../displays/sensitivity_global_timeseries.png}
\end{center}
\captext{Todos los métodos muestran tendencia positiva consistente desde $\sim$2008. La señal no depende de la elección metodológica.}
\end{frame}

%% ── Regional time series ─────────────────────────────────────────────────────
\begin{frame}{Series temporales regionales -- los 6 métodos}
\begin{center}
\fullimg{../displays/sensitivity_regional_timeseries.png}
\end{center}
\captext{El W. Pacific Warm Pool muestra la mayor sensibilidad al método. La Bahía de Bengala tiene el ciclo estacional más pronunciado.}
\end{frame}

%% ═══════════════════════════════════════════════════════════════════════════════
\sectionframe{5. Climatología Estacional}

%% ── Seasonal 4x6 ─────────────────────────────────────────────────────────────
\begin{frame}{BLD por estación: 4 estaciones $\times$ 6 métodos}
\begin{center}
\fullimg{../displays/monthly_clim_seasonal_comparison.png}
\end{center}
\captext{Pico de BLD en Sep--Oct en la mayoría de los métodos. El hemisferio sur muestra la señal estacional opuesta.}
\end{frame}

%% ── Clim Diff 0.2 ────────────────────────────────────────────────────────────
\begin{frame}{Climatología mensual -- Diferencia $\Delta T = 0.2$°C}
\begin{center}
\fullimg{../displays/monthly_clim_difference_02.png}
\end{center}
\captext{12 meses, 1993--2024. Escala compartida con el método de gradiente para comparación directa.}
\end{frame}

%% ── Clim Grad -0.025 ─────────────────────────────────────────────────────────
\begin{frame}{Climatología mensual -- Gradiente $-0.025$°C/m}
\begin{center}
\fullimg{../displays/monthly_clim_gradient_025.png}
\end{center}
\captext{12 meses, 1993--2024. Misma escala que la diapositiva anterior (vmax = 145.3 m).}
\end{frame}

%% ═══════════════════════════════════════════════════════════════════════════════
\sectionframe{6. Comparación LR 0.25° vs HR 0.083°}

%% ── Annual mean maps ─────────────────────────────────────────────────────────
\begin{frame}{Mapas anuales: LR, HR y diferencia (2024)}
\begin{center}
\fullimg{../displays/compare_lr_hr_annual_mean.png}
\end{center}
\captext{Media anual 2024. La diferencia LR--HR es pequeña en términos absolutos. Mayor desacuerdo donde hay alta variabilidad local.}
\end{frame}

%% ── Zoom regions ─────────────────────────────────────────────────────────────
\begin{frame}{Zoom en regiones clave -- diferencias LR vs HR}
\begin{center}
\fullimg{../displays/compare_lr_hr_zoom_regions.png}
\end{center}
\captext{La estructura espacial de la BLD se reproduce bien en ambas resoluciones.}
\end{frame}

%% ── Monthly stats ────────────────────────────────────────────────────────────
\begin{frame}{Series temporales mensuales de BLD: LR vs HR (2024)}
\begin{center}
\fullimg{../displays/compare_lr_hr_monthly_stats.png}
\end{center}
\captext{BLD media global mensual (arriba), diferencia HR--LR (abajo izq.) y cobertura de detección (abajo der.).}
\end{frame}

%% ── T/S scatter ──────────────────────────────────────────────────────────────
\begin{frame}{Acuerdo en perfiles T/S: scatter LR vs HR}
\begin{center}
\fullimg{../displays/scatter_profiles_lr_hr.png}
\end{center}
\captext{Columnas: Amazon Plume, Bahía de Bengala, W. Pacific Warm Pool. Filas: T (°C) y S (PSU). Cada punto = un par (mes, profundidad); color = profundidad.}
\end{frame}

%% ── BLD scatter ──────────────────────────────────────────────────────────────
\begin{frame}{Scatter global MLD/ILD/BLD: LR vs HR por cuenca (2024)}
\begin{center}
\fullimg{../displays/scatter_bld_lr_hr.png}
\end{center}
\captext{Mejor acuerdo en el Índico (R=0.88). Mayor dispersión en el Océano Austral (R=0.57). Sesgo MLD global: +3.0 m.}
\end{frame}

%% ── Resolvability Amazon Plume ────────────────────────────────────────────────
\begin{frame}{Resolvabilidad de diferencias ILD -- Amazon Plume}
\begin{center}
\fullimg{../displays/resolvability_amazon_plume.png}
\end{center}
\captext{Ratio = $\Delta$ILD / espaciado máximo de grilla. Ratio > 1 indica diferencia físicamente resolvable (no cuantización de grilla).}
\end{frame}

%% ── Resolvability Bay of Bengal ──────────────────────────────────────────────
\begin{frame}{Resolvabilidad de diferencias ILD -- Bahía de Bengala}
\begin{center}
\fullimg{../displays/resolvability_bay_of_bengal.png}
\end{center}
\captext{Muchas diferencias ILD entre LR y HR son cuantización de grilla, no señal física real.}
\end{frame}

%% ── Profiles with uncertainty Amazon ─────────────────────────────────────────
\begin{frame}{Perfiles LR vs HR con incertidumbre de grilla -- Amazon Plume}
\begin{center}
\fullimg{../displays/compare_profiles_amazon_plume_gradient_-0.1_with_uncertainty.png}
\end{center}
\captext{Bandas grises = $\pm\tfrac{1}{2}$ espaciado de grilla. Cuando se solapan, la diferencia ILD entre LR y HR no es resolvable.}
\end{frame}

%% ── Profiles with uncertainty Bay of Bengal ──────────────────────────────────
\begin{frame}{Perfiles LR vs HR con incertidumbre de grilla -- Bahía de Bengala}
\begin{center}
\fullimg{../displays/compare_profiles_bay_of_bengal_gradient_-0.1_with_uncertainty.png}
\end{center}
\captext{Bahía de Bengala: mayor acuerdo en meses de post-monzón. Diferencias resolvables principalmente en abr--may y ago.}
\end{frame}

%% ═══════════════════════════════════════════════════════════════════════════════
\sectionframe{7. Tendencias 1993--2024}

%% ── Trend table ──────────────────────────────────────────────────────────────
\begin{frame}{Tendencia global de BLD: los 6 métodos}
\begin{center}
\begin{tabular}{lrrrr}
\toprule
\textbf{Método} & \textbf{Tendencia} & \textbf{R} & \textbf{p-valor} & \textbf{Sig.} \\
                & (m/década)         &            &                  &               \\
\midrule
Gradiente $-0.015$°C/m        & +0.80 & 0.618 & 0.0002   & ** \\
\rowcolor{iceblue!60}
Gradiente $-0.025$°C/m        & +1.26 & 0.709 & $<$0.0001 & ** \\
Gradiente $-0.1$°C/m          & +1.75 & 0.705 & $<$0.0001 & ** \\
\rowcolor{iceblue!60}
Diff $\Delta T = 0.2$°C       & +0.43 & 0.278 & 0.123    & -- \\
Diff $\Delta T = 0.5$°C       & +1.08 & 0.521 & 0.002    & ** \\
Diff $\Delta T = 0.8$°C       & +1.31 & 0.574 & 0.001    & ** \\
\bottomrule
\end{tabular}
\end{center}
\vspace{4pt}
\captext{5 de 6 métodos: tendencia positiva significativa $\to$ señal física robusta, no artefacto metodológico.\\Resaltados en azul: métodos de referencia usados en las siguientes diapositivas.}
\end{frame}

%% ── Trend map Diff 0.2 all ────────────────────────────────────────────────────
\begin{frame}{Tendencia BLD -- Diferencia $\Delta T = 0.2$°C (todas las áreas)}
\begin{center}
\fullimg{../displays/bld_trend_map_difference_02_all.png}
\end{center}
\captext{Tendencia por píxel 1993--2024 (m/década). Escala: $\pm$12 m/década.}
\end{frame}

%% ── Trend map Diff 0.2 sig ────────────────────────────────────────────────────
\begin{frame}{Tendencia BLD -- Diferencia $\Delta T = 0.2$°C (solo p < 0.05)}
\begin{center}
\fullimg{../displays/bld_trend_map_difference_02_sig.png}
\end{center}
\captext{Solo píxeles con tendencia estadísticamente significativa (p < 0.05). El 19.9\% del océano; 90\% de ellos positivos.}
\end{frame}

%% ── Trend map Grad -0.025 all ─────────────────────────────────────────────────
\begin{frame}{Tendencia BLD -- Gradiente $-0.025$°C/m (todas las áreas)}
\begin{center}
\fullimg{../displays/bld_trend_map_gradient_025_all.png}
\end{center}
\captext{Tendencia por píxel 1993--2024 (m/década). Escala: $\pm$8 m/década.}
\end{frame}

%% ── Trend map Grad -0.025 sig ─────────────────────────────────────────────────
\begin{frame}{Tendencia BLD -- Gradiente $-0.025$°C/m (solo p < 0.05)}
\begin{center}
\fullimg{../displays/bld_trend_map_gradient_025_sig.png}
\end{center}
\captext{Solo píxeles con tendencia estadísticamente significativa (p < 0.05).}
\end{frame}

%% ── Regional trends ──────────────────────────────────────────────────────────
\begin{frame}{Tendencias regionales -- Gradiente $-0.025$°C/m (1993--2024)}
\begin{center}
\fullimg{../displays/bld_trend_regions.png}
\end{center}
\captext{Mayor tendencia: W. Pacific Warm Pool (+6.9 m/déc.). Señal más limpia: Índico Tropical (R=0.77). Bahía de Bengala: no significativa.}
\end{frame}

%% ── Trend ensemble ───────────────────────────────────────────────────────────
\begin{frame}{Spread del trend entre los 6 métodos}
\begin{center}
\fullimg{../displays/bld_trend_ensemble.png}
\end{center}
\captext{La incertidumbre en la magnitud del trend es $\sim2\times$ según el método. La dirección (positiva) es consistente en todos.}
\end{frame}

%% ═══════════════════════════════════════════════════════════════════════════════
\sectionframe{8. Variabilidad Interanual y Huella de ENSO}

%% ── Variability maps ─────────────────────────────────────────────────────────
\begin{frame}{Mapas de variabilidad temporal (1993--2024)}
\begin{center}
\fullimg{../displays/bld_variability_maps.png}
\end{center}
\captext{La variabilidad estacional (amplitud $\sim$29 m) domina sobre la interanual (std $\sim$6 m). Ratio $\approx5:1$ globalmente.}
\end{frame}

%% ── Anomaly maps ─────────────────────────────────────────────────────────────
\begin{frame}{Anomalías anuales de BLD (1993--2024)}
\begin{center}
\fullimg{../displays/bld_anomaly_maps.png}
\end{center}
\captext{Anomalías respecto a la media 1993--2024. Años extremos: 1998 (El Niño fuerte) y 2022 (La Niña 2021--23).}
\end{frame}

%% ── ENSO composite ───────────────────────────────────────────────────────────
\begin{frame}{Compuesto ENSO: anomalía de BLD en años Niño y Niña}
\begin{center}
\fullimg{../displays/enso_composite_gradient_025.png}
\end{center}
\captext{Método: gradiente $-0.025$°C/m. La Niña aumenta la BLD en el Pacífico occidental; El Niño la reduce. Der: diferencia La Niña -- El Niño.}
\end{frame}

%% ═══════════════════════════════════════════════════════════════════════════════
\sectionframe{9. Síntesis}

%% ── Key findings ─────────────────────────────────────────────────────────────
\begin{frame}{Hallazgos principales}
\begin{itemize}
  \item Las diferencias LR vs HR en perfiles T/S son mínimas (R > 0.997 para T, R > 0.978 para S)
  \item La grilla vertical LR es más fina que HR debajo de 30 m -- la resolución vertical no explica las diferencias de BLD
  \item La mayoría de las diferencias ILD entre LR y HR son cuantización de grilla, no señal física
  \item La BLD global aumentó +0.8 a +1.75 m/década (1993--2024), significativo en 5/6 métodos
  \item La señal más fuerte: W. Pacific Warm Pool (+6.9 m/década) e Índico Tropical (+2.4 m/década)
  \item La variabilidad estacional ($\sim$29 m de amplitud) supera $5\times$ a la interanual ($\sim$6 m)
  \item Huella de ENSO clara: mínimo en 1998 (El Niño), máximo en 2022 (La Niña)
\end{itemize}
\end{frame}

%% ── Next steps ───────────────────────────────────────────────────────────────
\begin{frame}{Próximos pasos}
\begin{itemize}
  \item Extender el análisis HR a años previos (disponible 1993--2024)
  \item Vincular la BLD con eventos de Marine Heatwaves (MHW)
  \item Analizar la conexión BLD--ENSO en profundidad (índices ONI, PDO)
  \item Comparar con datos observacionales (Argo floats)
  \item Definir el método óptimo para publicación basado en sensibilidad y resolvabilidad
\end{itemize}
\end{frame}

\end{document}
